% !TEX program = xelatex
\documentclass[letterpaper,12pt]{article}
\usepackage{kotex}
\usepackage[margin=0.75in]{geometry}
\usepackage{hyperref}

% Compact itemize spacing for one-page layout
\setlength{\leftmargini}{12pt}
\setlength{\leftmarginii}{12pt}
\makeatletter
\renewcommand{\@listI}{%
    \leftmargin\leftmargini
    \topsep 4pt \parsep 0pt \itemsep 5pt \partopsep 0pt
}
\let\@listi\@listI
\renewcommand{\@listii}{%
    \leftmargin\leftmarginii
    \labelwidth\leftmarginii
    \advance\labelwidth-\labelsep
    \topsep 2pt \parsep 0pt \itemsep 2pt \partopsep 0pt
}
\makeatother

% Disable paragraph indent
\setlength{\parindent}{0pt}
\linespread{1.12}\selectfont

% Custom section command for clean academic layout and readable spacing
\newcommand{\cvsection}[1]{%
    \vspace{15pt}%
    {\Large\bfseries #1}%
    \vspace{3pt}\hrule\vspace{6pt}%
}

% Entry header inside a section: bold title, right-aligned date
\newcommand{\cventry}[2]{%
    \vspace{9pt}%
    \textbf{#1} \hfill #2\par\vspace{2pt}%
}

% Sub-project header (one level below \cventry)
\newcommand{\cvsub}[1]{%
    \vspace{8pt}%
    \hspace*{10pt}\textbf{\textit{#1}}\par\vspace{2pt}%
}

\begin{document}

\begin{center}
    {\Huge \bfseries 이현중} \\ \vspace{6pt}
    서울과학기술대학교 전자IT미디어공학과 4학년 (2027년 2월 졸업 예정) \\ \vspace{4pt}
    \href{mailto:21101495@seoultech.ac.kr}{\texttt{21101495@seoultech.ac.kr}} | \href{https://github.com/bolyee}{\texttt{github.com/bolyee}}
\end{center}

\vspace{8pt}

음향을 공학적으로 탐구하고 구현하겠다는 목표를 가지고 전자공학과에 진학했습니다. 학부 과정에서 신호처리, 머신러닝, 회로 등 공학적 기반을 다지는 것에 그치지 않고, 다각적인 음향 실무를 현장에서 체득해 왔습니다. 음향 전문기업(A49,대표이사 이신렬 음향공학 박사)에서 음향기기 평가 프로그램 및 다양한 제품 개발 등 음향 실무를 익혔습니다. 또한 믹싱·마스터링 외주 작업과 단독 프로듀싱 음원 발매 등의 활동도 해왔습니다. 소리의 예술적인 부분, 객관적인 평가 그리고 공학적 기술을 모두 아우르는 음향공학자가 되고싶습니다.

% Section 1: Research Interests
\cvsection{Research Interests (관심 연구 분야)}
\begin{itemize}
    \item \textbf{ Sound Production \& Reproduction :} AI를 활용한 음악·오디오 제작 및 후처리, 사운드 합성·생성, 음색 변환과 제어
    \item \textbf{Auditory Perception :} 청각 심리학(Psychoacoustics), 청각 기반 음질 평가
    \item \textbf{Acoustic Signal Processing :} 음향 신호 분석, AI 생성 음원 판별, 오디오 이벤트 인식
    \item \textbf{Audio Measurement \& Electroacoustics :} 음향 기기 성능 측정 절차 설계 및 데이터 해석(사용장비: Audio Precision, NTI audio, Klippel)
\end{itemize}

% Section 2: Experience
\cvsection{Experience}

\cventry{서울과학기술대학교 전자공학과 신호처리 연구실 --- 학부 연구생}{2025.09 -- 현재}

\cvsub{AI 생성 음악 판별 모델 개발 (AI-Generated Music Detection)}
\begin{itemize}
    \item \textbf{동기:} AI 생성 음악이 많아지면서 이에 대한 사람들의 거부감이 확대되는 것을 보았습니다. 이에 따라 자동 필터링 시스템의 필요성을 확인하였습니다. 기존 모델의 경우 최신 모델로 생성된 AI 음악을 탐지하지 못한다는 한계점을 발견하였습니다.
    \item \textbf{수행 내용:} 사전학습된 모델을 직접 사용해보았고 최신 데이터 학습을 위한 파인튜닝을 진행하였습니다. 연구실 GPU 서버 활용을 위해 Docker를 익히고 의존성 문제를 해결하며, 학습 및 분류 과정을 경험해보았습니다.
    \item \textbf{한계 분석:} 공개 데이터셋 부족 등의 이유로 중단하였습니다. 향후 다시 연구해보고싶은 주제입니다.
\end{itemize}

\cvsub{급수음 기반 실시간 컵 수위 추정 (Pouring Sound Water Level Estimation)}
\begin{itemize}
    \item \textbf{동기:} 정수기 사용시 사용자가 출수 과정을 계속 지켜봐야 하는 불편함을 개선하고자 하였습니다.
    \item \textbf{수행 내용:} 급수음을 마이크로 수집하고, 실시간으로 컵의 수위를 추정하는 모델을 구현하였습니다. 급수를 제어하는 파이프라인은 ESP32로 구성하였습니다.
    \item \textbf{한계 분석:} 기존 데이터셋을 사용하며 컵의 정확한 수위를 예측하기 위해서는 컵을 사전 등록해야하는 제약이 존재하였습니다. 동일한 제약이라면 초음파 방식, 무게 센서 방식, 용량 사전 입력이 더 효율적이라 판단하여 주제를 전환하였습니다.
\end{itemize}

\cvsub{레퍼런스 곡 기반 자동 믹싱 시스템 (Reference-Based Auto-Mixing System)}
\begin{itemize}
    \item \textbf{동기:} DDSP(Differentiable DSP) 논문의 ``모델이 사전에 정의된 노브를 돌린다''는 관점을 실제 믹싱 워크플로우에 적용시켜 레퍼런스 곡 기반으로 믹싱을 자동화해보고 싶었습니다.
    \item \textbf{설계:} 레퍼런스 곡과 입력 간 차이를 손실 함수로 정의하여 오디오 이펙트 파라미터를 최적화하는 파이프라인을 구현하였습니다.
    \item \textbf{후속 계획:} 믹싱에서 중요한 요소인 Saturation으로 확장하여, 믹싱자동화시스템에 추가 및 새츄레이션과 지각과의 관계를 연구하고 싶습니다.
\end{itemize}

\cventry{A49(음향기기 연구개발회사, 대표이사: 이신렬 음향공학박사, 방학때마다 출근하여 연구를 진행) (\href{https://a49.co.kr/}{\texttt{https://a49.co.kr/}}) --- 연구보조 인턴 사원}{2024 -- 현재}
\begin{itemize}
    \item \textbf{노인성 난청 대상 청력보조 스피커} 개발에 참여하였습니다.  
    \item \textbf{해경 방송 시스템} 개발에 참여하였습니다. 하강풍으로 인해 방송소리가 들리지 않는 문제를 해결하기 위한 시스템입니다. CATT-Acoustic을 이용하여 간단한 음향 방사 특성과 음압 시뮬레이션을 진행하였습니다.
    \item \textbf{음향기기 객관 평가 시스템 구축 지원} (테크브로(TechBro) 홈페이지)(\href{https://techbros.co.kr/}{\texttt{https://techbros.co.kr/}}) 마이크로폰 객관적 평가 방법 개발 참여하였습니다.
    \item \textbf{음향 주관 평가 프로그램 개발}을 하였습니다. 귀마개(이어플러그) 음질과 감쇄성능 평가프로그램과 MUSHRA 기반 주관적 음질 평가 프로그램을 개발하였습니다.
\end{itemize}

\cventry{작곡 / 믹싱 · 마스터링 --- Music \& Audio Engineering}{2016 -- 현재}
\begin{itemize}
    \item \textbf{시작:} 중학생 시기 전자음악을 접한 뒤 FL Studio Mobile로 작곡을 시작하였습니다. ``프로의 사운드와 내 사운드는 왜 다른가''라는 질문이 현재 연구 관심의 출발점입니다.
    \item \textbf{학습:} FL Studio / Logic Pro 기반 작·편곡을 하였습니다. 군 복무 중 모바일 환경에서 믹싱·마스터링 이론을 학습하였으며, 전역 후 믹싱 수업을 수강하였습니다.
    \item \textbf{활동:} 작곡 동아리에서 믹싱 세미나(EQ, Compressor, 공간계 이펙터 원리 및 활용)를 진행하였습니다. 힙합 음원 믹싱 외주를 다수 수행하였습니다
    \item \textbf{음원 발매:} 2025년 작곡 동아리 앨범 프로젝트에서 작사 · 작곡 · 믹싱 · 마스터링을 단독으로 수행한 음원을 발매하였습니다. 교내 축제 무대에서 자작곡을 라이브로 연주하였습니다
    \item \textbf{활용 도구:} FL Studio, Logic Pro 및 다양한 플러그인 (Waves, iZotope 등)
\end{itemize}

% Section 3: Education
\cvsection{Education}
\textbf{서울과학기술대학교} \hfill 서울, 대한민국 \\
전자IT미디어공학과 학사 과정 (2021.03 -- 현재, 2027년 2월 졸업 예정) \\
\textit{학점:} 전공평점: 3.96/4.5 전체평점: 3.93/4.5 \\
\textit{주요 이수 과목:} 신호 및 시스템 (A+), 전력전자 (A+), 회로이론 (A+), 디지털신호처리, 공업수학 \\
\textit{외부 교육 이수:} 온라인 공개 강좌(MOOC) 기반 음향학(Acoustics) 교육 과정 수료(김양한 교수님), 믹싱/마스터링 교육 이수

\end{document}
